\documentclass[11pt]{amsart}
\usepackage{amsmath,amssymb,amsthm}
\usepackage[margin=1.1in]{geometry}
\usepackage{booktabs}
\usepackage[colorlinks=true,linkcolor=blue,citecolor=blue]{hyperref}

\makeatletter
\def\@settitle{\begin{center}%
  \baselineskip14\p@\relax
  \bfseries
  \@title
  \end{center}%
}
\def\@setauthors{%
  \begingroup
  \def\thanks{\protect\thanks@warning}%
  \trivlist
  \centering\footnotesize \@topsep30\p@\relax
  \advance\@topsep by -\baselineskip
  \item\relax
  \author@andify\authors
  \def\\{\protect\linebreak}%
  \authors
  \endtrivlist
  \endgroup
}
\let\ps@headings\ps@plain
\makeatother

\newtheorem{theorem}{Theorem}[section]
\newtheorem{proposition}[theorem]{Proposition}
\newtheorem{lemma}[theorem]{Lemma}
\newtheorem{corollary}[theorem]{Corollary}
\newtheorem{construction}[theorem]{Construction}
\newtheorem{problem}[theorem]{Problem}
\theoremstyle{definition}
\newtheorem{definition}[theorem]{Definition}
\newtheorem{remark}[theorem]{Remark}
\newtheorem{example}[theorem]{Example}

\newcommand{\C}{\mathbb{C}}
\newcommand{\Z}{\mathbb{Z}}
\newcommand{\F}{\mathbb{F}}
\newcommand{\ket}[1]{|#1\rangle}
\newcommand{\bra}[1]{\langle #1|}
\newcommand{\ip}[2]{\langle #1,#2\rangle}
\newcommand{\QLS}{\mathrm{QLS}}
\newcommand{\QLC}{\mathrm{QLC}}

\title[The Maximal Cardinality of Quantum Latin Cubes: A Complete Solution]
{The Maximal Cardinality of Quantum Latin Cubes:\\ A Complete Solution}

\author{Qingyu Zhang}
\address{School of Physics and Electronic Engineering, Jiangsu University\\
Zhenjiang 212013, Jiangsu, China}
\email{zhangqingyu2026@126.com}

\date{June 14, 2026}

\begin{document}

\begin{abstract}
% Plan: Define QLC and state the complete existence theorem, its constructions, and its structural consequences.
A quantum Latin cube (QLC) of order $v$ is a $v\times v\times v$ array of
unit vectors in $\C^v$ such that every line is an orthonormal basis.  Its
cardinality is the number of projectively distinct entries, at most $v^3$.
We prove that this maximum is attainable if and only if $v=1$ or $v\ge4$,
closing the previously open prime orders and orders $pq$ with $p,q$ prime
and $(p,q)\ne(2,2)$.  For every prime $p\ge5$, an explicit tri-clock
construction reduces maximality to linear independence of four functions
over $\F_p$.  Composite orders are handled by twisted layers over flat
quantum Latin squares.  We also characterise all layered QLCs in terms of
cyclic simple spectrum and flatness of the base square.  Consequently,
layered cardinalities are divisible by $v$, while the prime tri-clock family
has exact spectrum $\{p,p^2,p^3\}$.
\end{abstract}

\maketitle
\pagestyle{plain}
\makeatletter
\def\@oddhead{}%
\def\@evenhead{}%
\makeatother

\section{Introduction}\label{sec:intro}

A \emph{quantum Latin square} of order $v$, denoted $\QLS(v)$, is a $v\times
v$ array whose entries are unit column vectors of the Hilbert space
$\mathcal H_v=\C^v$ such that the $v$ entries of every row and of every
column form an orthonormal basis.  Quantum Latin squares were introduced by
Musto and Vicary~\cite{MV2016} and have since found applications to unitary
error bases, mutually unbiased bases, quantum Sudoku and absolutely maximally
entangled states~\cite{Musto2017,NP2021,Paczos2021,Goyeneche2018}.  Two unit
vectors $\ket u,\ket v$ are identified when $\ket u=e^{i\theta}\ket v$ for
some real $\theta$.  Following Paczos, Wierzbi\'nski,
Rajchel-Mield\'zio\'c, Burchardt and \.Zyczkowski~\cite{Paczos2021}, the
\emph{cardinality} $c$ of a $\QLS(v)$ is the number of distinct vectors
appearing in the array; clearly $v\le c\le v^2$.  After intensive recent
activity
\cite{Paczos2021,ZhangCao2026,ZZLC2026,ZZTS2025,ZWJ2026,ZhangJi2026} the
existence problem for squares of maximal cardinality is completely solved: a
$\QLS(v)$ with cardinality $v^2$ exists if and only if $v=1$ or
$v\ge4$~\cite{Paczos2021,ZhangCao2026,ZZTS2025}.

Quantum Latin \emph{cubes} were introduced by Goyeneche, Raissi, Di Martino
and \.Zyczkowski~\cite{Goyeneche2018} in connection with absolutely
maximally entangled states.  We use the definition of
\cite[Def.~1.2]{ZZLC2026}.

\begin{definition}\label{def:qlc}
A \emph{quantum Latin cube} of order $v$, denoted $\QLC(v)$, is a $v\times
v\times v$ array $C=(\ket{c_{i,j,k}})_{i,j,k\in\Z_v}$ of unit vectors of
$\C^v$ such that each \emph{line} of $C$ --- the $v$ entries obtained by
fixing two of the three indices --- forms an orthonormal basis of $\C^v$.
The \emph{cardinality} $c(C)$ is the number of distinct entries of $C$ up to
global phase; clearly $v\le c(C)\le v^3$.
\end{definition}

Replacing each symbol of a classical Latin cube by the corresponding
computational basis vector yields a $\QLC(v)$ of minimal cardinality $v$.
At the opposite extreme, Zhang, Zhang, Lv and Cao proved the following.

\begin{theorem}[{\cite[Thm.~1.3]{ZZLC2026}}]\label{thm:zzlc}
There exists a $\QLC(v)$ with maximal cardinality $v^3$ for all integers
$v\ge4$, except when $v$ is a prime or $v=pq$ with $p,q$ primes and
$(p,q)\ne(2,2)$.
\end{theorem}

Thus the existence of quantum Latin cubes of maximal cardinality remained
open for all primes $v\ge5$ and for all orders $pq$ with $p,q$ prime and
$(p,q)\ne(2,2)$, the smallest open orders being $v=5,6,7,9,10$.  In this paper we
resolve all remaining orders and obtain a complete answer, which exactly
parallels the dichotomy known for squares.

It is important that this is a cube endpoint problem rather than another
square-spectrum theorem.  The recent square results cited above settle the
maximal endpoint for $\QLS(v)$ and, for several infinite families, long or
complete intervals of square cardinalities.  They do not construct
three-dimensional arrays with $v^3$ projectively distinct entries in the
prime and $pq$ cases left open by Theorem~\ref{thm:zzlc}.  Conversely, our
main theorem determines only the maximal endpoint for cubes; it does not
claim a dense spectrum for $\QLC(v)$, nor does it classify orthogonal
families of quantum Latin squares or cubes.

\begin{theorem}\label{thm:main}
A quantum Latin cube of order $v$ with maximal cardinality $v^3$ exists if
and only if $v=1$ or $v\ge4$.
\end{theorem}

The proof has three components, each with a different role: a
rigidity theorem for layered cubes, an explicit construction for prime
orders, and a twisted-product construction for composite orders.

\emph{Prime orders: the tri-clock construction.}  For a prime $p$ and
permutations $\pi,g$ of $\Z_p$, define a cube whose entry at position
$(r,j,k)$ is the unit vector
\[
C_{r,j,k}\;=\;\frac1{\sqrt p}\,
\bigl(\,\omega^{\,rs+j\pi(s)+kg(s)}\,\bigr)_{s\in\Z_p},
\qquad \omega=e^{2\pi i/p}.
\]
Orthonormality along the three axes is immediate from the bijectivity of
$s\mapsto s$, $\pi$ and $g$, and we show
(Theorem~\ref{thm:triclock}) that $C$ has maximal cardinality $p^3$ if and
only if the four functions $1,\,s,\,\pi(s),\,g(s)$ are linearly independent
over the field $\F_p$.  More generally, its cardinality is $p^r$, where $r$
is the rank of $s,\pi(s),g(s)$ modulo the constant functions; this determines
the exact tri-clock spectrum $\{p,p^2,p^3\}$ for every $p\ge5$.  A single
explicit pair --- $\pi$ the transposition
of $p-2$ and $p-1$, and $g$ the $3$-cycle $0\mapsto1\mapsto2\mapsto0$ ---
satisfies this condition for \emph{every} prime $p\ge5$
(Lemma~\ref{lem:pair}), so the construction is completely explicit and its
verification is elementary.

\emph{Composite orders: twisted layers.}  For composite $v=mn\ge6$ we start
from the product construction of Zhang and Cao~\cite{ZhangCao2026,ZZLC2026},
which assembles a $\QLS(mn)$ of maximal cardinality out of row-quantum Latin
rectangles, and we choose the rectangles inside \emph{flat} (Fourier-type)
bases, so that every entry of the resulting square is unbiased to the
computational basis.  Stacking the $v$ layers $V^kQ$, $k\in\Z_v$, for a
suitable diagonal unitary $V$ then produces a $\QLC(v)$, and a combinatorial
choice of the eigenvalue assignment of $V$ forces all $v^3$ entries to be
distinct whenever one $2\times2$ double difference is a unit in $\Z_v$
(Theorem~\ref{thm:composite}).  Thus the composite-order construction is a
family, not a single example: for each factorisation $v=mn$ with $m,n\ge2$
and $v\ge6$ there are at least $v\varphi(v)(v-4)!$ admissible eigenvalue
assignments (Corollary~\ref{cor:many-composite}; inequivalence is not
asserted).  The layered architecture is not an ansatz
chosen for convenience: we prove a rigidity theorem
(Theorem~\ref{thm:rigidity}) showing that an array of the form
$(V^kq_{ij})_{i,j,k}$ is a quantum Latin cube \emph{precisely when} the
spectrum of $V$ is a rotated full set of $v$-th roots of unity and the layer
$Q=(q_{ij})$ is a quantum Latin square all of whose entries are unbiased to
the eigenbasis of $V$.  In particular the tri-clock cubes are exactly of
this form: their first layers are, up to relabelling, the
maximal-cardinality squares constructed by Zang, Zheng, Tian and
Shan~\cite[Thm.~3.1]{ZZTS2025} from phased Fourier matrices
(Corollary~\ref{cor:flatqls} and Remark~\ref{rem:zzts}).  Those squares are
unbiased to a fixed basis, and the rigidity theorem identifies this
flatness as precisely the property that makes a maximal-cardinality square
liftable to a cube --- a structural explanation of why this family, rather
than any other known family of maximal-cardinality squares, is the right
base for cube constructions.

The same framework gives two obstructions that delimit the reach of these
methods.  Every layered cube has cardinality divisible by its order
(Corollary~\ref{cor:layered-divisibility}), and a prime-order tri-clock cube
has cardinality $p^r$, where $r$ is the rank of three functions modulo the
constant functions (Theorem~\ref{thm:triclock}).  All three ranks occur, so
the tri-clock family has spectrum exactly $\{p,p^2,p^3\}$ for $p\ge5$
(Corollary~\ref{cor:triclock-spectrum}).  Thus neither construction can by
itself account for an arbitrary middle-spectrum value.

Finally, the same rigidity viewpoint explains why order four is genuinely
special: we show by a parity argument over $\F_2$ that no pair of
permutations of $\Z_4$ can satisfy the tri-clock independence condition
(Proposition~\ref{prop:four}), so the cube of \cite[Ex.~4.1]{ZZLC2026} cannot
be obtained by the tri-clock recipe; and we prove that quantum Latin cubes
of orders $2$ and $3$ are themselves necessarily classical, so that their
cardinality spectra are $\{2\}$ and $\{3\}$ (Theorem~\ref{thm:small}).
The key step for order $3$ is the elementary fact that a $3\times3$
unitary matrix with zero diagonal must be monomial
(Lemma~\ref{lem:monomial}).

As a supplementary application of the same projective-cardinality viewpoint,
we study quantum Latin squares of order five.  Combining exact constructions
with finite support--orthogonality certificates gives
$\operatorname{Spec}_{\QLS}(5)\cap[5,12]=\{5,7,12\}$
(Theorem~\ref{thm:qls5-low}).  Besides sharpening the small-order picture,
this calculation supplies a reproducible obstruction calculus for attacking
the still-open middle part of the square spectrum.  This square-spectrum
result is logically independent of the proof of Theorem~\ref{thm:main}.

The small-order square spectrum is changing quickly.  The dense-spectrum
theorems of Zhang, Wang and Ji~\cite{ZWJ2026} and Zhang and
Ji~\cite{ZhangJi2026} establish all possible cardinalities for the infinite
families of orders $4m$ and $6m$ with $m\ge2$, respectively, while Xu gives
explicit $\QLS(6)$ of cardinalities $13$, $15$ and
$17$~\cite{Xu2026}.  Thus the frontier for squares is no longer merely the
two endpoints; it is a mixture of dense high-divisibility results, isolated
small-order constructions, and finite obstruction problems.  This is the
context in which the certified order-five interval below should be read.

We separate throughout three levels of evidence.  The maximal-cardinality
cube theorem is proved by explicit formulae and finite algebraic checks.  The
order-five low-spectrum theorem is computer-assisted: its nonexistence
claims depend on a finite exhaustive enumeration of labelled
support--orthogonality types, after which every terminal type is eliminated
by one of the elementary linear-algebra configurations in
Lemma~\ref{lem:qls5-obstructions}.  Numerical searches mentioned in remarks
and problems are used only to guide conjectures; they are not evidence for
any theorem.

The paper is organised as follows.  Section~\ref{sec:prelim} fixes notation
and recalls the results we use.  Section~\ref{sec:layers} develops the
twisted-layer framework and the rigidity theorem.
Section~\ref{sec:composite} treats composite orders and
Section~\ref{sec:prime} prime orders; Theorem~\ref{thm:main} is assembled at
the end of Section~\ref{sec:prime}.  Section~\ref{sec:remarks} contains the
order-four obstruction, the small orders, the certified low spectrum of
$\QLS(5)$, and open problems on cardinality spectra.

\section{Preliminaries}\label{sec:prelim}

Throughout, $v\ge2$ is an integer, $\omega_v=e^{2\pi i/v}$, and
$\{\ket0,\dots,\ket{v-1}\}$ is the computational basis of $\C^v$.  Indices
of arrays are taken in $\Z_v$ and index arithmetic is modular.  For a unit
vector $\ket q$ and an orthonormal basis $B=\{\ket{b_0},\dots,\ket{b_{v-1}}\}$
we say $\ket q$ is \emph{flat} (or \emph{unbiased}) with respect to $B$ if
$|\ip{b_t}{q}|^2=1/v$ for all $t$.  The \emph{Fourier matrix} is
$F_v=\frac1{\sqrt v}\bigl(\omega_v^{st}\bigr)_{s,t\in\Z_v}$; its columns
form an orthonormal basis of flat vectors.

If $Q$ is a $\QLS(v)$ (or a $\QLC(v)$) and $U\in\mathrm U(v)$, the entrywise
image $UQ$ is again a $\QLS(v)$ (resp.\ $\QLC(v)$) with the same
cardinality, since unitaries preserve inner products and projective
distinctness.

We will use the following known facts.

\begin{lemma}[{\cite{Paczos2021}}]\label{lem:small-qls}
Every $\QLS(2)$ has cardinality $2$ and every $\QLS(3)$ has cardinality
$3$; equivalently, quantum Latin squares of orders $2$ and $3$ are
classical up to a global unitary.
\end{lemma}

\begin{lemma}[{\cite[Lem.~2.6]{ZhangCao2026}}]\label{lem:tensor}
Let $\ket a,\ket b\in\C^n$ and $\ket c,\ket d\in\C^m$ be unit vectors.  Then
$\ket a\otimes\ket c$ and $\ket b\otimes\ket d$ are equal up to phase if and
only if $\ket a,\ket b$ are equal up to phase and $\ket c,\ket d$ are equal
up to phase.
\end{lemma}

An \emph{$m\times n$ row-quantum Latin rectangle in $\C^n$} is an $m\times
n$ array of unit vectors of $\C^n$ each of whose rows is an orthonormal
basis of $\C^n$; its cardinality is the number of distinct entries.

\begin{lemma}[{\cite[Thm.~3.1]{ZhangCao2026}}; see also
{\cite[Lem.~2.2]{ZZLC2026}}]\label{lem:product}
Let $U=(\ket{u_{i,k}})$ be an $m\times n$ row-quantum Latin rectangle in
$\C^n$ with cardinality $c_1$ and $V=(\ket{v_{j,l}})$ an $n\times m$
row-quantum Latin rectangle in $\C^m$ with cardinality $c_2$.  Let $W$ be
the array of order $mn$ partitioned into an $m\times n$ grid of $n\times m$
blocks, whose entry in position $(k,l)$ of block $(i,j)$ is
\[
\ket{w_{i,j,k,l}}\;=\;\ket{u_{i,\,j+k}}\otimes\ket{v_{j,\,i+l}},
\qquad j+k \bmod n,\quad i+l \bmod m .
\]
Then $W$ is a $\QLS(mn)$ with cardinality $c_1c_2$.
\end{lemma}

We also record for context that quantum Latin squares of maximal
cardinality exist in every order $v\ge4$~\cite{ZZTS2025}.  Our
constructions below are self-contained and do not rely on this; the layer
square of the tri-clock cube does, however, turn out to coincide with the
construction of \cite[Thm.~3.1]{ZZTS2025} --- see Remark~\ref{rem:zzts}.

\section{Twisted layers and a rigidity theorem}\label{sec:layers}

The constructions of this paper produce a cube by stacking the iterates of
a fixed unitary applied to a single square.  In this section we determine
exactly when this is possible.

\begin{definition}
Let $V\in\mathrm U(v)$ and let $Q=(\ket{q_{i,j}})_{i,j\in\Z_v}$ be an array
of unit vectors of $\C^v$.  The \emph{$V$-layered cube over $Q$} is the
array $C=C(V,Q)$ with entries
\[
\ket{c_{i,j,k}}\;=\;V^{k}\ket{q_{i,j}},\qquad i,j,k\in\Z_v .
\]
\end{definition}

The lines of $C$ in the first two directions are the images under the
unitary $V^k$ of the rows and columns of $Q$, so they are orthonormal bases
if and only if $Q$ is a $\QLS(v)$.  The lines in the third direction --- the
\emph{orbit lines} $\{V^k\ket q\}_{k\in\Z_v}$ --- are governed by the
following two lemmas.

\begin{lemma}\label{lem:orbit}
For a unit vector $\ket q\in\C^v$ and $V\in\mathrm U(v)$, the orbit
$\{V^k\ket q:\ k=0,1,\dots,v-1\}$ is an orthonormal basis of $\C^v$ if and
only if $\ip{q}{V^k q}=0$ for $k=1,\dots,v-1$.
\end{lemma}

\begin{proof}
$\ip{V^aq}{V^bq}=\ip{q}{V^{\,b-a}q}$ and, for $0\le a<b\le v-1$, the
exponent $b-a$ ranges over $1,\dots,v-1$.
\end{proof}

\begin{theorem}[Rigidity]\label{thm:rigidity}
Let $V\in\mathrm U(v)$ and suppose some orthonormal basis
$\{\ket{q_1},\dots,\ket{q_v}\}$ of $\C^v$ has the property that every orbit
$\{V^k\ket{q_j}\}_{k\in\Z_v}$ is an orthonormal basis.  Then:
\begin{enumerate}
\item[\textup{(i)}] $V$ has $v$ distinct eigenvalues, and they form a
rotated full set of $v$-th roots of unity: there is a unimodular $\mu$ and
an eigenbasis $B=\{\ket{b_t}\}_{t\in\Z_v}$ with
$V\ket{b_t}=\mu\,\omega_v^{t}\ket{b_t}$;
\item[\textup{(ii)}] \emph{every} unit vector $\ket q$ whose orbit
$\{V^k\ket q\}_k$ is an orthonormal basis is flat with respect to $B$.
\end{enumerate}
\end{theorem}

\begin{proof}
Let $\lambda_1,\dots,\lambda_r$ be the distinct eigenvalues of $V$ with
spectral projections $P_1,\dots,P_r$, and for a unit vector $\ket q$ put
$w_u(q)=\|P_u\ket q\|^2$, so that $w_u(q)\ge0$, $\sum_u w_u(q)=1$ and
\begin{equation}\label{eq:moments}
\ip{q}{V^kq}\;=\;\sum_{u=1}^{r} w_u(q)\,\lambda_u^{\,k},
\qquad k\in\Z .
\end{equation}
Fix $j$ and abbreviate $w_u=w_u(q_j)$.  By Lemma~\ref{lem:orbit},
$\sum_u w_u\lambda_u^k=0$ for $k=1,\dots,v-1$.

Suppose $r\le v-1$.  Then the equations with $k=1,\dots,r$ are available,
and they say that the vector $(w_u\lambda_u)_u$ is annihilated by the
Vandermonde matrix $(\lambda_u^{\,k-1})_{k=1,\dots,r;\ u=1,\dots,r}$, which
is invertible because the $\lambda_u$ are distinct.  Hence
$w_u\lambda_u=0$, so $w_u=0$ for all $u$, contradicting $\sum_u w_u=1$.
Therefore $r=v$: all eigenvalues of $V$ are simple, and the eigenbasis
$B=\{\ket{b_t}\}$ is determined up to phases; write
$w_t(q)=|\ip{b_t}{q}|^2$.

Now the full system $\sum_t w_t\lambda_t^{\,k}=\delta_{k,0}$
($k=0,1,\dots,v-1$) is a square Vandermonde system, so it has a
\emph{unique} solution $w^*=(w^*_t)_t$, and $w(q)=w^*$ for every unit
vector $\ket q$ whose orbit is an orthonormal basis --- in particular for
each $\ket{q_j}$.  Since $\{\ket{q_j}\}_j$ is an orthonormal basis,
Parseval gives, for each $t$,
\[
v\,w^*_t\;=\;\sum_{j=1}^{v}|\ip{b_t}{q_j}|^2\;=\;\|\ket{b_t}\|^2\;=\;1,
\]
so $w^*_t=1/v$: every such $\ket q$ is flat with respect to $B$, proving
(ii).  Finally, \eqref{eq:moments} with $w=w^*$ yields
$\sum_t\lambda_t^{\,k}=0$ for $k=1,\dots,v-1$, i.e.\ the first $v-1$ power
sums of the $\lambda_t$ vanish; by Newton's identities all elementary
symmetric functions $e_1,\dots,e_{v-1}$ vanish as well, so the
characteristic polynomial of $V$ is $z^v-(-1)^{v+1}e_v$ and the
$\lambda_t$ are the $v$-th roots of a single unimodular number.  This is
(i).
\end{proof}

\begin{theorem}[Characterisation of layered cubes]\label{thm:layered}
$C(V,Q)$ is a $\QLC(v)$ if and only if
\begin{enumerate}
\item[\textup{(a)}] $Q$ is a $\QLS(v)$;
\item[\textup{(b)}] the spectrum of $V$ is a rotated full set of $v$-th
roots of unity, $\{\mu\,\omega_v^t:\ t\in\Z_v\}$, with all eigenvalues
simple;
\item[\textup{(c)}] every entry of $Q$ is flat with respect to the
eigenbasis of $V$.
\end{enumerate}
\end{theorem}

\begin{proof}
($\Leftarrow$)  Lines in the first two directions are unitary images of
rows and columns of $Q$.  For an orbit line, (b) and (c) give, for
$k\not\equiv0$,
\[
\ip{q}{V^kq}=\sum_t |\ip{b_t}{q}|^2(\mu\omega_v^t)^k
=\frac{\mu^k}{v}\sum_{t\in\Z_v}\omega_v^{tk}=0 ,
\]
and Lemma~\ref{lem:orbit} applies.
($\Rightarrow$)  The layer $k=0$ of a $\QLC$ is a $\QLS$, giving (a).  The
first row of $Q$ is an orthonormal basis each of whose orbits is a line of
$C$, hence an orthonormal basis; Theorem~\ref{thm:rigidity} gives (b), and
(ii) of that theorem applies to every entry of $Q$ (each entry's orbit is a
line of $C$), giving (c).
\end{proof}

\begin{lemma}[Cardinality of layered cubes]\label{lem:card}
In the situation of Theorem~\ref{thm:layered}, write $\ket q\sim_V\ket{q'}$
if $\ket{q'}$ is proportional to $V^k\ket q$ for some $k$.  Then each
projective $V$-orbit $\{[V^k\ket q]:k\in\Z_v\}$ has exactly $v$ elements,
and
\[
c\bigl(C(V,Q)\bigr)\;=\;v\cdot\#\bigl(\{\text{entries of }Q\}/\!\sim_V\bigr).
\]
In particular $c(C)=v^3$ if and only if $c(Q)=v^2$ and no two distinct
entries of $Q$ are $\sim_V$-equivalent.
\end{lemma}

\begin{proof}
If $V^k\ket q=e^{i\theta}\ket q$ with $0<k<v$, then comparing coordinates
in the eigenbasis and using that a flat vector has no zero coordinate,
$\mu^k\omega_v^{tk}=e^{i\theta}$ for all $t$, whence
$\omega_v^{(t-t')k}=1$ for all $t,t'$, i.e.\ $v\mid k$, a contradiction.
So all orbits are free, and two orbits are either disjoint or equal as
projective sets.  The entries of $C$ are precisely the union of the orbits
of the entries of $Q$, and the count follows.
\end{proof}

\begin{corollary}[Layered divisibility obstruction]
\label{cor:layered-divisibility}
The cardinality of every layered $\QLC(v)$ is divisible by $v$.  Hence a
$\QLC(v)$ whose cardinality is not divisible by $v$ admits no representation
$C(V,Q)$, even after permuting the three coordinate directions.
\end{corollary}

\begin{proof}
The divisibility follows from Lemma~\ref{lem:card}.  Permuting the coordinate
directions preserves cardinality, so the same obstruction applies to a
layered representation in any direction.
\end{proof}

\section{Composite orders}\label{sec:composite}

\subsection{Flat Fourier bases}

\begin{lemma}\label{lem:flatbases}
For $n\ge2$ and $i\in\Z_{\ge0}$ define unit vectors
$\ket{x^{(i)}_a}\in\C^n$, $a\in\Z_n$, by
\[
\bigl(x^{(i)}_a\bigr)_s\;=\;\frac1{\sqrt n}\;
e^{\,i\sqrt2\,\cdot i\,\delta_{s,0}}\;\omega_n^{\,sa},
\qquad s\in\Z_n,
\]
that is, $B_i:=\{\ket{x^{(i)}_a}:a\in\Z_n\}$ consists of the columns of
$D_iF_n$ where $D_i=\operatorname{diag}(e^{i\sqrt2\,i},1,\dots,1)$.  Then
each $B_i$ is an orthonormal basis all of whose vectors are flat with
respect to the computational basis, and for $i\ne i'$ the bases $B_i$ and
$B_{i'}$ have no projective vector in common.
\end{lemma}

\begin{proof}
$D_iF_n$ is unitary with all entries of modulus $n^{-1/2}$, giving the
first claim.  Suppose $\ket{x^{(i)}_a}=e^{i\theta}\ket{x^{(i')}_{a'}}$.
Comparing coordinates,
\begin{equation}\label{eq:coincide}
\sqrt2\,(i-i')\,\delta_{s,0}+\frac{2\pi s}{n}(a-a')\;\equiv\;\theta
\pmod{2\pi}\qquad\text{for all } s\in\Z_n .
\end{equation}
If $n\ge3$, the cases $s=1,2$ give $\tfrac{2\pi}{n}(a-a')\equiv0$, so
$a=a'$ and then $\theta\equiv0$, and $s=0$ gives
$\sqrt2\,(i-i')\in2\pi\Z$; since $\pi$ is transcendental this forces
$i=i'$.  If $n=2$, the case $s=1$ gives $\theta\equiv\pi(a-a')$, so
$\theta\in\pi\Z$, and $s=0$ gives $\sqrt2\,(i-i')\in\pi\Z$, again forcing
$i=i'$ and then $a=a'$.
\end{proof}

\subsection{A flat square of maximal cardinality}

Let $v=mn$ with $m,n\ge2$.  Apply Lemma~\ref{lem:flatbases} twice: in
dimension $n$ it yields bases $B_0,\dots,B_{m-1}$ of $\C^n$, giving the
$m\times n$ row-quantum Latin rectangle $U=(\ket{u_{i,a}})$ with
$\ket{u_{i,a}}=\ket{x^{(i)}_a}$; in dimension $m$ it yields bases
$B'_0,\dots,B'_{n-1}$ of $\C^m$, giving the $n\times m$ row-quantum Latin
rectangle $V=(\ket{v_{j,b}})$.  All $mn$ entries of $U$ are pairwise
projectively distinct, and likewise for $V$, so $c_1=c_2=mn$.  Let $W$ be
the $\QLS(mn)$ of Lemma~\ref{lem:product}; then $c(W)=(mn)^2$ is maximal.
Moreover:

\begin{enumerate}
\item every entry of $W$ equals $\ket{u_{i,a}}\otimes\ket{v_{j,b}}$ for
some $(i,a)\in\Z_m\times\Z_n$ and $(j,b)\in\Z_n\times\Z_m$, and conversely
every such tensor occurs as an entry (in block $(i,j)$, at inner position
$k=a-j$, $l=b-i$);
\item every entry of $W$ is flat with respect to the computational product
basis
\[
\{\ket s\otimes\ket t:\ (s,t)\in\Z_n\times\Z_m\},
\]
because each coordinate of $\ket{u_{i,a}}\otimes\ket{v_{j,b}}$ has modulus
$n^{-1/2}m^{-1/2}=v^{-1/2}$.
\end{enumerate}

\subsection{The twist}

\begin{lemma}\label{lem:assignment}
For $v=mn\ge6$ with $m,n\ge2$ there exists a bijection
$g:\Z_n\times\Z_m\to\Z_v$ with a \emph{unit double difference}, i.e.
\[
\Delta g:=g(1,1)-g(1,0)-g(0,1)+g(0,0)\in\Z_v^\times .
\]
In fact one may prescribe
\[
g(0,0)=0,\qquad g(1,0)=1,\qquad g(0,1)=2,\qquad g(1,1)=4 .
\]
Consequently, for each fixed factorisation $v=mn$, there are at least
$v\varphi(v)(v-4)!$ such bijections, where $\varphi$ is Euler's function.
\end{lemma}

\begin{proof}
The four cells exist since $m,n\ge2$, the four prescribed values are
distinct elements of $\Z_v$ since $v\ge5$, and the remaining $v-4$ values
$\{3,5,6,\dots,v-1\}$ may be assigned to the remaining cells arbitrarily.
For every such assignment $\Delta g=4-1-2+0=1$, a unit in $\Z_v$.

This gives $(v-4)!$ normalised bijections.  If $a\in\Z_v^\times$ and
$b\in\Z_v$, then $ag+b$ is again a bijection and
\[
\Delta(ag+b)=a\,\Delta g\in\Z_v^\times .
\]
The resulting $v\varphi(v)(v-4)!$ bijections are distinct: from $h=ag+b$
one recovers $b=h(0,0)$ and $a=h(1,0)-h(0,0)$, and then recovers $g$.
\end{proof}

\begin{theorem}\label{thm:composite}
Let $v=mn\ge6$ with $m,n\ge2$, and let
$g:\Z_n\times\Z_m\to\Z_v$ be any bijection with unit double difference
$\Delta g\in\Z_v^\times$.  Then the twisted-product construction below gives
a $\QLC(v)$ with maximal cardinality $v^3$.  In particular every composite
integer $v\ge6$ admits a $\QLC(v)$ with maximal cardinality.
\end{theorem}

\begin{proof}
Let $W$ be the flat $\QLS(v)$ of maximal cardinality constructed above, and
let
\[
V_g\;=\;\sum_{(s,t)\in\Z_n\times\Z_m}\omega_v^{\,g(s,t)}\,
(\ket s\otimes\ket t)(\bra s\otimes\bra t) .
\]
The spectrum of $V_g$ is the full set of $v$-th roots of unity with simple
eigenvalues, and every entry of $W$ is flat with respect to its eigenbasis
(the computational product basis), so $C=C(V_g,W)$ is a $\QLC(v)$ by
Theorem~\ref{thm:layered}.

It remains to verify, via Lemma~\ref{lem:card}, that distinct entries of
$W$ are $\sim_{V_g}$-inequivalent.  Suppose
\[
V_g^{\,\kappa}\bigl(\ket{u_{i,a}}\otimes\ket{v_{j,b}}\bigr)
=e^{i\theta}\,\ket{u_{i',a'}}\otimes\ket{v_{j',b'}}
\]
for some $\kappa\in\Z_v$.  All coordinates of both sides are nonzero, and
comparing their phases at the coordinate $(s,t)$ gives, for all
$(s,t)\in\Z_n\times\Z_m$,
\begin{equation}\label{eq:collision}
\frac{2\pi\kappa}{v}\,g(s,t)\;+\;A(s)\;+\;B(t)\;\equiv\;\theta
\pmod{2\pi},
\end{equation}
where
$A(s)=\sqrt2\,(i-i')\,\delta_{s,0}+\tfrac{2\pi s}{n}(a-a')$ and
$B(t)=\sqrt2\,(j-j')\,\delta_{t,0}+\tfrac{2\pi t}{m}(b-b')$ collect all
terms depending on $s$ alone, respectively on $t$ alone.  Evaluating
\eqref{eq:collision} at the four cells $(0,0),(1,0),(0,1),(1,1)$ and
taking the alternating sum annihilates $\theta$, $A$ and $B$ entirely and
leaves
\[
\frac{2\pi\kappa}{v}\,\bigl[g(1,1)-g(1,0)-g(0,1)+g(0,0)\bigr]
=\frac{2\pi\kappa\,\Delta g}{v}\;\equiv\;0\pmod{2\pi},
\]
whence $v\mid\kappa\Delta g$.  Since $\Delta g$ is a unit modulo $v$, this
gives $\kappa\equiv0$.  Then the two entries of $W$
coincide projectively, and since $c(W)=v^2$, Lemma~\ref{lem:tensor} and
Lemma~\ref{lem:flatbases} force $(i,a,j,b)=(i',a',j',b')$.  By
Lemma~\ref{lem:card}, $c(C)=v\cdot v^2=v^3$.
The existence assertion follows by choosing $g$ as in
Lemma~\ref{lem:assignment}.
\end{proof}

\begin{corollary}[Abundance of composite twists]\label{cor:many-composite}
For each factorisation $v=mn\ge6$ with $m,n\ge2$, the preceding construction
has at least $v\varphi(v)(v-4)!$ admissible eigenvalue assignments $g$ giving
maximal-cardinality layered $\QLC(v)$.  This count is by assignments; no
inequivalence of the resulting cubes is asserted.
\end{corollary}

\begin{proof}
Use the assignments of Lemma~\ref{lem:assignment} in
Theorem~\ref{thm:composite}.
\end{proof}

\begin{corollary}\label{cor:composite}
For every composite integer $v\ge4$ there exists a $\QLC(v)$ with maximal
cardinality.
\end{corollary}

\begin{proof}
Combine Theorem~\ref{thm:composite} with the explicit $\QLC(4)$ of maximal
cardinality of \cite[Ex.~4.1]{ZZLC2026}.
\end{proof}

\begin{remark}
The hypothesis $v\ge6$ in Lemma~\ref{lem:assignment} is sharp for this
corner-twist method.  If $g:\Z_2\times\Z_2\to\Z_4$ is a bijection, then,
modulo $2$,
\[
\Delta g\equiv g(1,1)+g(1,0)+g(0,1)+g(0,0)
\equiv0+1+0+1\equiv0,
\]
so $\Delta g$ is even and cannot be a unit in $\Z_4$.  The tri-clock method
has its own order-four obstruction, proved in Proposition~\ref{prop:four}.
\end{remark}

\section{Prime orders: the tri-clock construction}\label{sec:prime}

Throughout this section $p\ge5$ is prime and $\omega=\omega_p$.  We regard
functions $\Z_p\to\Z_p$ as vectors of the $p$-dimensional space
$\F_p^{\Z_p}$ over the field $\F_p$; the four functions of interest are the
constant $\mathbf1$, the identity $\iota(s)=s$, and two permutations
$\pi,g$ of $\Z_p$.

\begin{theorem}\label{thm:triclock}
Let $p$ be prime and let $\pi,g$ be permutations of $\Z_p$.  Define the
cube $C=(\ket{c_{r,j,k}})_{r,j,k\in\Z_p}$ by
\[
\ket{c_{r,j,k}}\;=\;\frac1{\sqrt p}
\Bigl(\omega^{\,rs+j\pi(s)+kg(s)}\Bigr)_{s\in\Z_p}\;\in\;\C^p .
\]
Then $C$ is a $\QLC(p)$.  If
\[
r_0=\dim_{\F_p}\operatorname{span}
\bigl\{[\iota],[\pi],[g]\bigr\}
\quad\text{in}\quad
\F_p^{\Z_p}/\langle\mathbf1\rangle ,
\]
then
\[
c(C)=p^{r_0}.
\]
In particular $r_0\in\{1,2,3\}$, so a tri-clock cube has cardinality
$p$, $p^2$ or $p^3$; it has maximal cardinality if and only if
$\{\mathbf1,\iota,\pi,g\}$ is linearly independent over $\F_p$.
\end{theorem}

\begin{proof}
\emph{Lines.}  For a line in the first direction (fix $j,k$, vary $r$),
\[
\ip{c_{r,j,k}}{c_{r',j,k}}
=\frac1p\sum_{s\in\Z_p}\omega^{(r'-r)s}
=\delta_{r,r'} ,
\]
since $\sum_{s\in\Z_p}\omega^{cs}=0$ for $c\not\equiv0$.  For the second
direction the same computation applies after the substitution
$u=\pi(s)$, which is a bijection of $\Z_p$:
$\ip{c_{r,j,k}}{c_{r,j',k}}=\frac1p\sum_u\omega^{(j'-j)u}=\delta_{j,j'}$;
likewise for the third direction using $g$.  Hence $C$ is a $\QLC(p)$.

\emph{Cardinality.}  The projective class of the entry at $(r,j,k)$ is
determined by the coset of
$r\iota+j\pi+kg$ modulo the constant functions.  More explicitly, two
entries coincide projectively,
$\ket{c_{r,j,k}}=e^{i\theta}\ket{c_{r',j',k'}}$, if and only if the
function
\[
s\;\longmapsto\;(r-r')\,s+(j-j')\,\pi(s)+(k-k')\,g(s)\pmod p
\]
is constant on $\Z_p$, i.e.\ if and only if
$(r-r')\,\iota+(j-j')\,\pi+(k-k')\,g\in\F_p\mathbf1$.
Thus the projective classes of the entries are in bijection with the image
of the linear map
\[
\F_p^3\longrightarrow \F_p^{\Z_p}/\langle\mathbf1\rangle,\qquad
(r,j,k)\longmapsto[r\iota+j\pi+kg].
\]
Its image has dimension $r_0$, proving $c(C)=p^{r_0}$.  Since $\iota$ is
non-constant, $1\le r_0\le3$.  If
$\{\mathbf1,\iota,\pi,g\}$ is linearly independent this forces
$r=r'$, $j=j'$, $k=k'$, so all $p^3$ entries are distinct.  Conversely, a
nontrivial relation $a\mathbf1+b\iota+c\pi+dg=0$ must have
$(b,c,d)\ne(0,0,0)$ (else $a=0$ as well), and then
$b\,s+c\,\pi(s)+d\,g(s)\equiv-a$ is constant, so the entries at positions
$(0,0,0)$ and $(b,c,d)$ coincide projectively and $c(C)<p^3$.
\end{proof}

\begin{lemma}\label{lem:pair}
For every prime $p\ge5$, let $\pi$ be the transposition exchanging $p-2$
and $p-1$ (and fixing all other points), and let $g$ be the $3$-cycle
$g(0)=1$, $g(1)=2$, $g(2)=0$ (fixing all other points).  Then
$\{\mathbf1,\iota,\pi,g\}$ is linearly independent over $\F_p$.
\end{lemma}

\begin{proof}
Suppose $a+bs+c\,\pi(s)+d\,g(s)\equiv0\pmod p$ for all $s\in\Z_p$.

First let $p\ge7$.  The set $S=\{3,4,\dots,p-3\}$ has $p-5\ge2$ elements,
and for $s\in S$ we have $\pi(s)=g(s)=s$, so $a+(b+c+d)s\equiv0$ at two or
more distinct values of $s$; hence $b+c+d\equiv0$ and $a\equiv0$.  At
$s=p-2$ and $s=p-1$ we have $g(s)=s$, so $(b+d)s+c\,\pi(s)\equiv0$;
subtracting the two instances gives $(b+d)-c\equiv0$, i.e.\ $c\equiv b+d$.
Combining with $b+c+d\equiv0$ yields $2c\equiv0$, so $c\equiv0$ (as $p$ is
odd) and $d\equiv-b$.  Finally $s=0$ gives $d\,g(0)=d\equiv0$, hence
$b\equiv0$.  Thus $a=b=c=d=0$.

For $p=5$ we have $\pi=(0,1,2,4,3)$ and $g=(1,2,0,3,4)$ (in one-line
notation), and the rows $s=0,1,2,3$ of the matrix
$\bigl(\mathbf1,\iota,\pi,g\bigr)$ form
\[
M=\begin{pmatrix}
1&0&0&1\\ 1&1&1&2\\ 1&2&2&0\\ 1&3&4&3
\end{pmatrix},
\qquad
\det M
=\det\!\begin{pmatrix}1&1&2\\2&2&0\\3&4&3\end{pmatrix}
-\det\!\begin{pmatrix}1&1&1\\1&2&2\\1&3&4\end{pmatrix}
=4-1=3\not\equiv0\pmod 5 ,
\]
so the four columns are linearly independent over $\F_5$.
\end{proof}

\begin{corollary}\label{cor:prime}
For every prime $p\ge5$ there exists an explicit $\QLC(p)$ with maximal
cardinality $p^3$.
\end{corollary}

\begin{corollary}[Exact spectrum of the prime tri-clock family]
\label{cor:triclock-spectrum}
For every prime $p\ge5$, the set of cardinalities attained by tri-clock
$\QLC(p)$ is exactly
\[
\{p,p^2,p^3\}.
\]
\end{corollary}

\begin{proof}
Theorem~\ref{thm:triclock} shows that no other cardinality can occur.  Taking
$\pi=g=\iota$ gives rank $r_0=1$ and cardinality $p$.  If $\pi$ is any
non-affine permutation and $g=\pi$, then $[\iota]$ and $[\pi]$ are linearly
independent modulo the constants, giving rank $r_0=2$ and cardinality $p^2$.
Such a permutation exists for $p\ge5$, for example the transposition in
Lemma~\ref{lem:pair}: an affine map fixing three points must be the identity,
whereas this transposition is not.  Finally, the pair in Lemma~\ref{lem:pair}
has rank $r_0=3$ and gives cardinality $p^3$.
\end{proof}

\begin{example}\label{ex:p5}
For $p=5$, with $\pi=(0,1,2,4,3)$ and $g=(1,2,0,3,4)$, the cube
\[
\ket{c_{r,j,k}}=\tfrac1{\sqrt5}
\bigl(\omega_5^{\,rs+j\pi(s)+kg(s)}\bigr)_{s\in\Z_5},
\qquad r,j,k\in\Z_5,
\]
is a $\QLC(5)$ with cardinality $125$.  All inner products are fifth-root
geometric sums and the verification is exact; a machine check (residuals
$\sim10^{-16}$, no projective coincidences among the $125$ entries) is
provided in the accompanying scripts, together with the analogous
verifications for $p=7,11,13,17$.
\end{example}

\begin{remark}\label{rem:layered-view}
The tri-clock cube is layered in the sense of Section~\ref{sec:layers}:
$C=C(V_g,Q)$ where $V_g=\sum_s\omega^{\,g(s)}\ket s\bra s$ and $Q$ is the
layer $k=0$,
\[
\ket{q_{r,j}}=\frac1{\sqrt p}\bigl(\omega^{\,rs+j\pi(s)}\bigr)_{s\in\Z_p}.
\]
All entries of $Q$ are flat with respect to the computational basis, in
accordance with Theorem~\ref{thm:layered}.  Note also that all $p^3$
entries of $C$ are flat: the cube is, as a whole, unbiased to one fixed
basis.
\end{remark}

\begin{corollary}[{cf.\ \cite[Thm.~3.1]{ZZTS2025}}]\label{cor:flatqls}
For every prime $p\ge5$ the square $Q$ of Remark~\ref{rem:layered-view}
(with $\pi$ any non-affine permutation of $\Z_p$, e.g.\ the transposition
of Lemma~\ref{lem:pair}) is a $\QLS(p)$ of maximal cardinality $p^2$ all of
whose entries are unbiased to the computational basis.
\end{corollary}

\begin{proof}
Rows and columns are orthonormal as in Theorem~\ref{thm:triclock}.  Two
entries coincide iff $(r-r')\iota+(j-j')\pi\in\F_p\mathbf1$; if
$j\ne j'$ this exhibits $\pi$ as an affine function
$\pi(s)=(j-j')^{-1}\bigl(c-(r-r')s\bigr)$, and if $j=j'$, $r\ne r'$ it
makes $(r-r')\iota$ constant; both are impossible.  Non-affine
permutations exist for $p\ge5$ because $p!>p(p-1)$.
\end{proof}

\begin{remark}\label{rem:zzts}
Up to relabelling, the square $Q$ is the family of maximal-cardinality
quantum Latin squares constructed, for \emph{all} orders $v\ge4$, by Zang,
Zheng, Tian and Shan~\cite[Thm.~3.1]{ZZTS2025}, who phase the columns of
the Fourier matrix by the columns of a row-permuted Fourier matrix.  The
present formulation adds two points.  First, at prime orders the
non-affineness of the permutation $\pi$ is \emph{necessary} as well as
sufficient for maximal cardinality (the proof of
Corollary~\ref{cor:flatqls} is an equivalence).  Second, all entries of
these squares are unbiased to the computational basis, and by
Theorem~\ref{thm:layered} this flatness is exactly the property required
of the layer of a quantum Latin cube --- which is why this family, rather
than any other known family of maximal-cardinality squares, lifts to the
cubes of Corollary~\ref{cor:prime}.  By Musto's correspondence between
quantum Latin squares and mutually unbiased bases~\cite{Musto2017}, flat
squares are also a natural source of MUB-type structures; we do not pursue
this here.
\end{remark}

We can now prove the main theorem.

\begin{proof}[Proof of Theorem~\textup{\ref{thm:main}}]
For $v=1$ the unique $\QLC(1)$ has cardinality $1=v^3$.  For $v\in\{2,3\}$
maximal cardinality is impossible by Theorem~\ref{thm:small} below.
For $v\ge4$ composite, apply Corollary~\ref{cor:composite}; for $v\ge5$
prime, apply Corollary~\ref{cor:prime}.
\end{proof}

\section{Small orders, the order-four obstruction, and open problems}
\label{sec:remarks}

\begin{lemma}\label{lem:monomial}
A $3\times3$ unitary matrix with zero diagonal is monomial: it has exactly
one nonzero entry in every row and every column.
\end{lemma}

\begin{proof}
Let the columns be $b_0=(0,\beta,\gamma)^{\top}$,
$b_1=(\delta,0,\varepsilon)^{\top}$, $b_2=(\zeta,\eta,0)^{\top}$.
Orthogonality of the three pairs of columns gives
$\bar\gamma\varepsilon=\bar\beta\eta=\bar\delta\zeta=0$.  If $\gamma=0$,
then $|\beta|=1$, so $\eta=0$, so $|\zeta|=1$, so $\delta=0$, and the
matrix is monomial.  If instead $\varepsilon=0$, then $|\delta|=1$, so
$\zeta=0$, so $|\eta|=1$, so $\beta=0$, and again the matrix is monomial.
\end{proof}

\begin{theorem}\label{thm:small}
Every $\QLC(2)$ has cardinality $2$, and every $\QLC(3)$ has cardinality
$3$: quantum Latin cubes of orders $2$ and $3$ are classical up to a
global unitary.  In particular the cardinality spectra of quantum Latin
cubes of orders $2$ and $3$ are $\{2\}$ and $\{3\}$, and maximal
cardinality $v^3$ is unattainable.
\end{theorem}

\begin{proof}
\emph{Order $2$.}  Each layer $\{\ket{c_{i,j,k}}\}_{i,j}$ (fixed $k$) is a
$\QLS(2)$, so by Lemma~\ref{lem:small-qls} it carries the two vectors of
an orthonormal basis of $\C^2$.  The file condition forces
$\ket{c_{i,j,1}}\perp\ket{c_{i,j,0}}$, and in $\C^2$ the orthogonal
complement of a unit vector is projectively unique; hence layer $1$
carries the same two projective vectors as layer $0$, and $c(C)=2$.

\emph{Order $3$.}  By Lemma~\ref{lem:small-qls}, layer $k$ is classical:
there are a Latin square $L_k$ of order $3$ and an orthonormal basis
\[
B_k=\{\ket{b^{(k)}_0},\ket{b^{(k)}_1},\ket{b^{(k)}_2}\}
\]
with
$\ket{c_{i,j,k}}\propto\ket{b^{(k)}_{L_k(i,j)}}$.  After a global unitary
we may assume $B_0=\{\ket{e_0},\ket{e_1},\ket{e_2}\}$ is the computational
basis.  We claim $B_1=B_0$ projectively.

For $y\in\{0,1,2\}$ let $N(y)=\{L_0(i,j):\ L_1(i,j)=y\}$.  The file
through a cell $(i,j)$ with $L_1(i,j)=y$ contains $\ket{c_{i,j,0}}$, so
$\ket{b^{(1)}_y}\perp\ket{e_x}$ for every $x\in N(y)$.

Suppose first that $|N(y_0)|=1$ for some $y_0$, say $N(y_0)=\{x_0\}$.
The three cells of the symbol $y_0$ in $L_1$ then carry all three
occurrences of $x_0$ in $L_0$, so $x_0\notin N(y)$ for $y\ne y_0$ and
$N(y_1),N(y_2)\subseteq\{x_1,x_2\}$.  If some $|N(y_1)|=2$, then each of
$x_1,x_2$ occurs at least once among the cells of $y_1$ and hence at most
twice among those of $y_2$, while $y_2$ has three cells; consequently
$N(y_2)=\{x_1,x_2\}$ as well.  But then $\ket{b^{(1)}_{y_1}}$ and
$\ket{b^{(1)}_{y_2}}$ are both orthogonal to $\ket{e_{x_1}}$ and
$\ket{e_{x_2}}$, hence both proportional to $\ket{e_{x_0}}$,
contradicting their orthogonality.  Therefore in this case
$|N(y)|=1$ for all $y$, and $y\mapsto N(y)$ is a bijection (two symbols of
$L_1$ cannot pair exclusively with the same symbol of $L_0$, which occurs
only three times).  The matrix of $B_1$ in computational coordinates then
has, after relabelling the columns, zero diagonal, so it is monomial by
Lemma~\ref{lem:monomial}: $B_1=B_0$ projectively.

Otherwise $|N(y)|\ge2$ for every $y$, and each $\ket{b^{(1)}_y}$, being
orthogonal to two distinct computational vectors, is proportional to the
third: again $B_1\subseteq B_0$, hence $B_1=B_0$ as both are bases.

The same argument gives $B_2=B_0$, so all $27$ entries are computational
up to phase and $c(C)=3$.
\end{proof}

Theorem~\ref{thm:small} is consistent with the numerical exploration of
the $\QLC(3)$ variety: $60$ out of $60$ random alternating-projection runs
converged, all of them to classical solutions of cardinality $3$ (script
\texttt{verify\_qlc3\_rigidity.py}).

\begin{proposition}[Order four obstruction]\label{prop:four}
There is no pair of permutations $\pi,g$ of $\Z_4$ such that
$s\mapsto \alpha s+\beta\pi(s)+\gamma g(s) \bmod4$ is non-constant for
every $(\alpha,\beta,\gamma)\in\Z_4^3\setminus\{(0,0,0)\}$.  Consequently
the tri-clock recipe of Theorem~\ref{thm:triclock} cannot produce a cube of
maximal cardinality at order $4$.
\end{proposition}

\begin{proof}
Let $\rho:\Z_4\to\F_2$ be reduction modulo $2$.  For any permutation $h$
of $\Z_4$ the vector $\rho\circ h\in\F_2^{\Z_4}\cong\F_2^4$ is
\emph{balanced} (two $0$s and two $1$s), since $h$ hits two even and two
odd values.  In the quotient space $\F_2^4/\langle\mathbf1\rangle$
(dimension $3$) the six balanced vectors fall into exactly three classes,
namely $[0101]$, $[0011]$ and $[0110]$, and these satisfy
$[0011]+[0110]=[0101]$; together with $0$ they form a $2$-dimensional
subspace.  The classes $[\rho\circ\iota]=[0101]$, $[\rho\circ\pi]$ and
$[\rho\circ g]$ therefore lie in a common $2$-dimensional space, so they
are linearly dependent over $\F_2$: there are
$\varepsilon_1,\varepsilon_2,\varepsilon_3\in\F_2$, not all zero, and
$\varepsilon_0\in\F_2$ with
$\varepsilon_1\rho(s)+\varepsilon_2\rho(\pi(s))+\varepsilon_3\rho(g(s))
+\varepsilon_0=0$ in $\F_2$ for all $s$.  Multiplying by $2$ and using
$2\rho(x)\equiv2x\pmod4$, we obtain
\[
2\varepsilon_1\,s+2\varepsilon_2\,\pi(s)+2\varepsilon_3\,g(s)
\;\equiv\;-2\varepsilon_0\pmod 4\qquad\text{for all } s ,
\]
a constant combination with
$(\alpha,\beta,\gamma)=(2\varepsilon_1,2\varepsilon_2,2\varepsilon_3)
\ne(0,0,0)$.
\end{proof}

\begin{remark}\label{rem:composite-triclock}
For composite $v$ the tri-clock cube of Theorem~\ref{thm:triclock} is still
a $\QLC(v)$ for any permutations $\pi,g$ of $\Z_v$ (the line computations
only use bijectivity).  Its cardinality is the size of the image of the
$\Z_v$-module homomorphism
\[
\Z_v^3\longrightarrow
\Z_v^{\Z_v}/\langle\mathbf1\rangle,\qquad
(\alpha,\beta,\gamma)\longmapsto
[\alpha\iota+\beta\pi+\gamma g].
\]
In particular its cardinality divides $v^3$.  It has maximal cardinality if
and only if
$\alpha\iota+\beta\pi+\gamma g$ is non-constant modulo $v$ for all
$(\alpha,\beta,\gamma)\ne(0,0,0)$ --- equivalently, if for every prime
$q\mid v$ the reductions $\{\mathbf1,\iota,\pi,g\}\bmod q$ are linearly
independent over $\F_q$ as functions on $\Z_v$.  Proposition~\ref{prop:four}
shows this fails for $v=4$; an exhaustive computer search confirms that no
pair of permutations of $\Z_4$ works.  For $v=6$ pairs do exist --- e.g.\
$\pi=(0,1,2,3,5,4)$, $g=(0,1,2,4,3,5)$ give a tri-clock $\QLC(6)$ of
cardinality $216$, verified by machine --- so the twisted-product
construction of Section~\ref{sec:composite} and the tri-clock construction
overlap at many composite orders.  A characterisation of the orders $v$
admitting tri-clock pairs would be of independent combinatorial interest.
\end{remark}

\begin{problem}\label{prob:composite-triclock}
Determine the possible image sizes of the module homomorphism in
Remark~\ref{rem:composite-triclock} as $\pi,g$ range over permutations of
$\Z_v$.  In particular, characterise the composite orders $v$ for which
there are permutations $\pi,g$ such that
$\alpha\iota+\beta\pi+\gamma g$ is non-constant modulo $v$ for every
nonzero $(\alpha,\beta,\gamma)\in\Z_v^3$.  Equivalently, by the Chinese
remainder theorem, characterise when the reductions of
$\{\mathbf1,\iota,\pi,g\}$ remain linearly independent over every prime
field $\F_q$ with $q\mid v$.  The obstruction above rules out $v=4$, while
the example in Remark~\ref{rem:composite-triclock} shows that $v=6$ is
already possible.
\end{problem}

\subsection{The certified low spectrum of \texorpdfstring{$\QLS(5)$}{QLS(5)}}
\label{sec:qls5-spectrum}

Write $\operatorname{Spec}_{\QLS}(n)$ for the set of cardinalities attained
by quantum Latin squares of order $n$.  The order-five case already displays
substantial gaps.  The following result combines the exact and
computer-assisted conclusions of the order-five investigation.  Here
\emph{certified} means that, after the first-row normalisation, the remaining
possibilities are reduced to a finite list of support--edge types and every
type is eliminated by an explicitly recorded obstruction from
Lemma~\ref{lem:qls5-obstructions}.  No floating-point feasibility test enters
the nonexistence proof.

\begin{lemma}[Support--orthogonality obstructions]\label{lem:qls5-obstructions}
Let $x_1,\dots,x_t$ be distinct non-computational projective vectors of
$\C^5$, with prescribed coordinate supports of size at least two and with
some prescribed orthogonalities.  Each of the following configurations is
impossible.
\begin{enumerate}
\item[\textup{(a)}] $q$ pairwise orthogonal vectors have their supports
contained in a coordinate subspace of dimension less than $q$.
\item[\textup{(b)}] Distinct $x,y,z$ lie in a common $2$-space and satisfy
$x\perp y\perp z$.
\item[\textup{(c)}] Orthogonal $x,y$ span a coordinate $2$-space $V$, while
$z\perp x,y$ has fewer than two support coordinates outside $V$.
\item[\textup{(d)}] A vector $x$ with $2$-element support is orthogonal to
$y$, and their supports meet in exactly one coordinate, without $y$
vanishing at that coordinate.
\item[\textup{(e)}] Distinct $u,w$ have the same $3$-element support $S$ and
two common orthogonal neighbours $a,b$ whose $2$-element supports lie in
$S$, meet in one coordinate, and together cover $S$.
\end{enumerate}
\end{lemma}

\begin{proof}
Part (a) is the dimension bound for pairwise orthogonal nonzero vectors.
For (b), the orthogonal complement of $y$ inside the common $2$-space is a
line containing both $x$ and $z$.  In (c), the projection of $z$ onto $V$
is orthogonal to the basis $\{x,y\}$, so $z$ is supported outside $V$ and
therefore degenerates to a computational vector.  For (d), if
$\operatorname{supp}(x)=\{i,j\}$ and the intersection is $\{i\}$, then
$0=\langle x,y\rangle=\overline{x_i}y_i$; both coordinates of the
non-computational vector $x$ are nonzero, hence $y_i=0$.  Finally, in (e)
the support conditions make $a,b$ linearly independent, and the orthogonal
complement of their span in $\C^S$ is a line containing both $u$ and $w$.
\end{proof}

\begin{theorem}[Computer-assisted certified low spectrum of order five]
\label{thm:qls5-low}
\[
\operatorname{Spec}_{\QLS}(5)\cap[5,12]=\{5,7,12\}.
\]
More precisely, cardinalities $5$, $7$ and $12$ are attained, whereas
$6$, $8$, $9$, $10$ and $11$ are not.
\end{theorem}

\begin{proof}
Cardinality $5$ is attained by every classical Latin square.  For
cardinality $7$, take the Latin square with rows
\[
01234,\quad 10342,\quad 23401,\quad 34120,\quad 42013.
\]
Its first two rows and first two columns contain an intercalate on symbols
$0,1$; rotating those two symbols by a generic angle in their coordinate
$2$-space introduces exactly two new projective vectors.  Construction
\ref{con:qls5-c12} gives cardinality $12$.  The exclusion of $6=5+1$ is the
general $n+1$ nonexistence lemma of~\cite[Lem.~3.1]{ZWJ2026}.

It remains to exclude $8,9,10,11$.  Normalize the first row to the
computational basis.  A square of cardinality $c=5+t$ then has $t$ new
projective classes on the remaining $20$ cells.  Its \emph{active graph} is
the bipartite graph on the four lower rows and five columns whose edges are
the cells carrying new classes.  Every complete labelling determines
coordinate supports for the new classes and a forced orthogonality graph.
The deterministic certificate has three layers: enumerate active graphs up
to row and column permutations; enumerate all admissible label patterns
with new labels introduced in canonical first-occurrence order; and
classify the resulting support--edge types by
Lemma~\ref{lem:qls5-obstructions}.  The exhaustive terminal statistics,
counted after symmetry reduction at the active-graph layer, are
\[
\begin{array}{ccrrr}
c & t & \text{graph types} & \text{label patterns} &
\text{support--edge types}\\
\hline
8  & 3 & 5   & 34\,560  & \text{directly classified}\\
9  & 4 & 14  & 1\,587   & 97\\
10 & 5 & 40  & 9\,010   & 1\,724\\
11 & 6 & 220 & 201\,297 & 46\,808
\end{array}
\]
and no type survives.  Thus none of $8,9,10,11$ is attainable.  The table is
not a numerical approximation: each row records a finite certificate whose
leaves are disposed of by Lemma~\ref{lem:qls5-obstructions}.
\end{proof}

\begin{remark}[Proof boundary]\label{rem:qls5-proof-boundary}
The exact constructions and Lemma~\ref{lem:qls5-obstructions} are ordinary
mathematical proofs.  The exclusions of cardinalities $8$--$11$ additionally
depend on the exhaustiveness and correctness of the enumeration programs.
The machine-readable outputs record zero unclassified terminal types; the
entry points and certificate files are listed in the computational note.
\end{remark}

\begin{construction}[An exact cardinality-$12$ square]\label{con:qls5-c12}
Let $e_0,\dots,e_4$ be the computational vectors and define the following
unnormalised integer vectors:
\[
\begin{aligned}
a&=(0,-1,0,2,1), & b&=(0,1,0,1,-1), & c&=(0,1,0,0,1),\\
d&=(0,-1,0,2,-5), & h&=(0,2,0,1,0), & f&=(0,1,0,-2,0),\\
g&=(0,1,0,0,-1). &&
\end{aligned}
\]
After normalising these seven vectors, the array
\[
Q_{12}=
\begin{pmatrix}
e_0 & e_1 & e_2 & e_3 & e_4\\
a   & e_0 & b   & c   & e_2\\
d   & e_2 & a   & e_0 & h\\
h   & e_4 & e_0 & e_2 & f\\
e_2 & e_3 & c   & g   & e_0
\end{pmatrix}
\]
is a $\QLS(5)$ of cardinality exactly $12$.  All row and column
orthogonality checks reduce to integer dot products before normalisation.
\end{construction}

\begin{remark}\label{rem:qls5-next}
The next unresolved low value is $13$.  A tangent-space calculation at
$Q_{12}$ gives projective tangent nullity $4$ after fixing the first row,
and the nullity remains $4$ under every tested single-cell split of a
repeated class; this is evidence of first-order rigidity, not a
nonexistence proof.  High-precision numerical candidates of cardinalities
$21$ and $24$ have also been found, while extensive stratified searches suggest
the conjectural full spectrum
$\operatorname{Spec}_{\QLS}(5)=\{5,7,12,21,24,25\}$.
At present we make no exact claim about $13,\dots,20,22,23$.
\end{remark}

\begin{problem}\label{prob:spectrum}
Determine the \emph{cardinality spectrum} of quantum Latin cubes: for which
$c\in[v,v^3]$ does a $\QLC(v)$ of cardinality $c$ exist?
Theorem~\ref{thm:small} settles the question for $v\le3$: the spectra are
$\{2\}$ and $\{3\}$.  For $v\ge4$ both endpoints are attained (the
classical cube, and Theorem~\ref{thm:main}), and Lemma~\ref{lem:card}
shows that every flat $\QLS(v)$ stratum of cardinality $c'$ with $N$
twist-orbit classes contributes the value $vN$ to the spectrum.  Exploratory
alternating projections on the $\QLC(5)$ variety (script
\path{qlc_experiments.py}) reached, besides the value $125$ of
Corollary~\ref{cor:prime}, the cardinalities $21$, $25$ and $45$.  These
are numerical candidates rather than theorems, but they indicate that the
spectrum of $\QLC(5)$ is nontrivial in its middle range.  The candidate
$21$ would necessarily be non-layered by
Corollary~\ref{cor:layered-divisibility}; and Theorem~\ref{thm:triclock}
shows that neither $21$ nor $45$ can arise from the tri-clock construction,
whose order-five cardinalities are restricted to $5$, $25$ and $125$.  The
corresponding spectrum problem for squares is an active subject
\cite{Paczos2021,ZWJ2026,ZhangJi2026}; Theorem~\ref{thm:qls5-low} determines
its order-five low-cardinality interval, but the middle range remains open.
Even the first attainable cardinality strictly above $v$, and whether long
intervals of cardinalities occur for any fixed $v$, are unknown in general.
\end{problem}

\begin{problem}\label{prob:flat-spectrum}
Corollary~\ref{cor:flatqls} produces, in every prime dimension $p\ge5$,
quantum Latin squares of maximal cardinality all of whose $p^2$ entries
are unbiased to a fixed basis.  Which cardinalities are achievable by such
\emph{flat} quantum Latin squares?  By Theorem~\ref{thm:layered} this
question, together with the orbit quotient in Lemma~\ref{lem:card}, governs
all layered quantum Latin cubes.  In $173$ converged
random least-squares searches over the flat $\QLS(5)$ variety we observed
only the two cardinalities $5$ and $25$, suggesting that the flat spectrum
may be much sparser than the unrestricted one.
\end{problem}

\begin{problem}\label{prob:qls5-middle}
Complete the order-five square spectrum.  The first exact target is the
value $13$, where the same support--edge certificate framework applies but
the number of dense active graphs becomes the bottleneck.  A second target
is to replace the numerical candidates of cardinalities $21$ and $24$ by exact
constructions, or to prove that they cannot be rationalised inside any
low-dimensional algebraic family related to $Q_{12}$.
\end{problem}

\begin{problem}\label{prob:small-square-spectra}
Determine the complete spectra of $\QLS(5)$ and $\QLS(6)$.  For order five
the certified beginning is Theorem~\ref{thm:qls5-low}.  For order six, the
general $n+1$ obstruction excludes $7$, the dense theorem of
\cite{ZhangJi2026} applies only from order $12$ onward, and the explicit
order-six constructions of~\cite{Xu2026} show that the small spectrum has
new isolated values beyond the classical and maximal endpoints.  A complete
classification of these two orders would test whether the sparse behaviour
of $\QLS(5)$ is exceptional or the first instance of a broader small-order
phenomenon.
\end{problem}

\subsection*{Concluding perspective}
Theorem~\ref{thm:main} closes the existence problem at the maximal endpoint:
apart from the rigid small orders $2$ and $3$, every order admits a quantum
Latin cube with all entries projectively distinct.  The proof also isolates
the structural requirement behind the layered method.  A square can serve
as a layer precisely when it is flat in an eigenbasis whose associated
unitary has cyclic simple spectrum; maximality then becomes an
orbit-separation problem.  This separates two tasks that were previously
intertwined and provides a framework for future spectrum constructions.
Corollary~\ref{cor:layered-divisibility} and Theorem~\ref{thm:triclock}
also turn this framework into an obstruction theory: layered cardinalities
are multiples of $v$, and the prime-order tri-clock spectrum is exactly
$\{p,p^2,p^3\}$.  On the constructive side,
Corollary~\ref{cor:many-composite} shows that the composite-order endpoint is
not isolated: every nontrivial factorisation $v=mn\ge6$ supplies a factorial
family of admissible twists, with at least $v\varphi(v)(v-4)!$ choices of
eigenvalue assignment.

The result does not determine non-maximal cube cardinalities, nor does it
classify maximal cubes up to equivalence.  Those problems require control of
orbit collisions rather than their complete avoidance.  The certified
order-five square calculation illustrates one possible complementary tool:
finite support--orthogonality obstructions can detect gaps that explicit
constructions alone do not reveal.

\subsection*{Computational note}
All constructions used in the proofs are exact.  The tri-clock entries are
roots of unity, the composite construction uses Fourier entries together
with the transcendental phase $e^{i\sqrt2\,r}$ only to separate projective
classes, and the cardinality-$12$ square uses normalised integer vectors.
The corresponding checks reduce to finite geometric sums, phase-comparison
identities, or integer dot products.

The computer-assisted nonexistence result for
$\operatorname{Spec}_{\QLS}(5)\cap[8,11]$ is finite and deterministic:
the scripts enumerate active graphs, label patterns and support--edge
types, then attach one of the obstructions of
Lemma~\ref{lem:qls5-obstructions} to each terminal type.  The numerical
least-squares and alternating-projection experiments mentioned in
Remarks~\ref{rem:qls5-next} and Problems~\ref{prob:spectrum},
\ref{prob:flat-spectrum} are exploratory only.

Machine-verification scripts for the cube constructions accompany this draft
in the main source directory and include
\path{verify_uniform_pair.py} (tri-clock cubes for
$p=5,7,11,13,17$), \path{verify_composite.py} (twisted-product cubes for
$v=6,9,10$), \path{triclock_experiments.py} (the exhaustive order-$4$
search, the order-$6$ pair, and the count of $8000$ valid pairs
$(\pi,g)$ in $S_5\times S_5$), and \path{verify_qlc3_rigidity.py} (the
order-$3$ rigidity check).  The order-five square-spectrum scripts and
machine-readable certificates are in the companion directory
\path{qls_spectrum/}.  Its principal entry points are:
\begin{itemize}
\item \path{c8_classification_certificate.py} and
\path{qls5_c9_deterministic_certificate.py};
\item \path{qls5_c10_deterministic_probe.py} and
\path{qls5_incremental_probe.py};
\item \path{qls_support_obstruction_general.py} and
\path{exact_qls5_c12_construction.py}.
\end{itemize}
The corresponding JSON outputs record zero unclassified terminal types.
The companion report \path{qls5_low_cardinality_report.tex} records the full
certificate statistics.

\begin{thebibliography}{10}

\bibitem{Goyeneche2018}
D.~Goyeneche, Z.~Raissi, S.~Di~Martino and K.~\.Zyczkowski,
\emph{Entanglement and quantum combinatorial designs},
Phys.\ Rev.\ A \textbf{97} (2018), 062326.

\bibitem{MV2016}
B.~Musto and J.~Vicary,
\emph{Quantum Latin squares and unitary error bases},
Quantum Inf.\ Comput.\ \textbf{16} (2016), 1318--1332.

\bibitem{Musto2017}
B.~Musto,
\emph{Constructing mutually unbiased bases from quantum Latin squares},
Electron.\ Proc.\ Theor.\ Comput.\ Sci.\ \textbf{236} (2017), 108--126.

\bibitem{NP2021}
I.~Nechita and J.~Pillet,
\emph{SudoQ --- a quantum variant of the popular game},
Quantum Inf.\ Comput.\ \textbf{21} (2021), 781--799.

\bibitem{Paczos2021}
J.~Paczos, M.~Wierzbi\'nski, G.~Rajchel-Mield\'zio\'c, A.~Burchardt and
K.~\.Zyczkowski,
\emph{Genuinely quantum solutions of the game Sudoku and their cardinality},
Phys.\ Rev.\ A \textbf{104} (2021), 042423.

\bibitem{ZhangCao2026}
Y.~Zhang and H.~Cao,
\emph{Quantum Latin squares with maximal cardinality},
Discrete Math.\ \textbf{349} (2026), 114863.

\bibitem{ZZLC2026}
Y.~Zhang, Y.~Zhang, M.~Lv and H.~Cao,
\emph{The maximal cardinality of quantum Latin squares and quantum Latin
cubes},
J.\ Combin.\ Des.\ \textbf{34} (2026), 169--183.

\bibitem{ZZTS2025}
Y.~Zang, M.~Zheng, Z.~Tian and X.~Shan,
\emph{On the cardinalities of quantum Latin squares},
arXiv:2508.01972 (2025).

\bibitem{ZWJ2026}
Y.~Zhang, X.~Wang and L.~Ji,
\emph{Quantum Latin squares with all possible cardinalities},
J.\ Combin.\ Des.\ \textbf{34} (2026), 378--387.

\bibitem{ZhangJi2026}
Y.~Zhang and L.~Ji,
\emph{Quantum Latin squares of order $6m$ with all possible cardinalities},
arXiv:2601.09132 (2026).

\bibitem{Xu2026}
Z.~Xu,
\emph{Three quantum Latin squares of order $6$ with cardinalities $13$, $15$
and $17$},
arXiv:2605.15540 (2026).

\end{thebibliography}

\end{document}
